\chapter{Transfer}
\label{chap:transfer}
Im Folgendem wird die in Kapitel~\ref{chap:methodik} beschriebene Methodik auf den MDV-Prozess angewendet, um den bislang nicht modellierten Ablauf transparent darzustellen. Zuerst wird die Erstaufnahme des MDV-Ablaufs beschrieben und daraufhin die Modellierung mit BPMN 2.0 in Signavio. Danach wird das  Modell in der Prozessarchitektur eingeordnet.

\section{Erstaufnahme des MDV-Ablaufs}
\label{sec:erstaufnahme}
Gemäß Abschnitt~\ref{sec:prozessaufnahme} wurden die Prozessschritte, die Verantwortlichen sowie die Eingangs- und Ausgangsgrößen des MDV-Prozesses erhoben. Damit das Modell für alle beteiligten Fachbereiche verständlich bleibt, wurde der Ablauf in Hauptschritte mit jeweils einer Tätigkeitsbeschreibung gegliedert. Da das Prozesswissen auf einzelne Fachexperten verteilt war und kein dokumentierter Ablauf vorlag, wurden drei Interviews mit Vertretern aus drei Fachbereichen geführt, die aufgrund ihrer Beteiligung unterschiedliche Abschnitte des MDV-Prozesses abdecken konnten. Eine Einschränkung der Aufnahme bestand darin, dass die Interviews zu Beginn ausschließlich die Perspektiven einzelner Fachexperten widerspiegelten. Eine Dokumentenanalyse wurde nicht durchgeführt, da die vorhandenen Unterlagen den tatsächlichen Ablauf nicht vollständig abbildeten. Auf einen gemeinsamen Workshop wurde auch verzichtet, um die unterschiedlichen Sichtweisen der Fachbereiche getrennt zu erfassen. 
\\
Als Ergebnis der Erstaufnahme wurden 13 Prozessschritte vom Kick-off bis zur Übermittlung des finalen MDV erfasst. Um die Ergebnisse abzusichern, wurde das  entwickelte Prozessmodell gemeinsam mit den direkt Beteiligten fachlich abgestimmt. Die in Kapitel~\ref{chap:problemkontext} angenommene Abfolge diente einer ersten Strukturierung der bestehenden Black Box und wird durch den nun aufgenommenen Ablauf abgelöst. Die Durchführung liegt beim Team Konformitäts-Testing, während übrige Fachbereiche konsultiert oder informiert werden. Die beteiligten Rollen sind in Tabelle \ref{tab:rollen} aufgelistet. Als Eingangsgrößen wurden u. a. der Software-Reifegrad und der Terminplan der Zertifizierung aufgenommen. Daneben wurden drei Entscheidungspunkte erfasst, bei denen die Vollständigkeit der Ergebnisse sowie die Bestätigung des Pre-MDV und des finalen MDV geprüft werden. Die Ausgangsgröße bildet der finale MDV, der schließelich im System eingetragen und an eine weitere Abteilung übermittelt wird. 

\section{Modellierung und Beschreibung des MDV-Prozesses in Signavio}
\label{sec:modellierung}
Der MDV-Prozess wurde auf Grundlage der Ergebnisse aus Abschnitt \ref{sec:erstaufnahme} in Signavio modelliert, da unternehmensweite Prozesse bei Mercedes-AMG auf dieser Plattform abgebildet werden \citeintern{signavio2026intranet}. Die Abbildung \ref{fig:MDV-Prozess} zeigt das Modell für den Markt Westeuropa, das die 13 Prozessschritte abbildet. Weil das ursprüngliche Modell zu breit für eine lesbare Darstellung auf einer Seite ist, wird eine vereinfachte Fassung gezeigt, ohne Rollen, Datenobjekte, Systeme sowie Tätigkeitsbeschreibungen. Die Elemente der Notation aus der Abbildung \ref{fig:BPMN} wurden für den vorliegenden Prozess verwendet. Hierbei wurden die Aktivitäten einem Pool für Mercedes-AMG und einer Lane für das Team Konformitäts-Testing zugeordnet, was den Prozess durchführt. Die übrigen Fachbereiche wurden als Rollen an den Aktivitäten in Tabelle \ref{tab:rollen} hinterlegt, da bis zu elf Rollen pro Aktivität die Darstellung als Lane unübersichtlich gemacht hätten.
\newline
Der modellierte Ablauf gliedert sich in drei Phasen. Der Prozess beginnt mit dem Startereignis \glqq AZST Eintrag Zert. Vorgang abgeschlossen\grqq, sobald der Zertifizierungsvorgang im Abgas-Zertifizierungs-Steuerungstool (\acs{AZST}) eingetragen wurde. Im anschließenden Kick-off werden die Fachbereiche konsultiert und liefern den Input, aus dem die MDV-Meilensteine für den Terminplan Zertifizierung entstehen. Daraufhin prüft ein Gateway, ob alle notwendigen Ergebnisse vorliegen, wobei fehlende Informationen zunächst eingeholt werden, bevor beide Pfade wieder zusammenlaufen. Anschließend werden die Vortestergebnisse konsolidiert, wobei die Ergebnisse des Ausrollprozesses über ein UND-Gateway einfließen, und die Messwerte bewertet. Die erforderlichen Informationen werden in der Aktivität \glqq Notwendige Daten zusammengeführt\grqq gebündelt.
\newline
In der zweiten Phase wird daraufhin auf Basis der vorherigen Phase der Pre-MDV ermittelt, dessen Output der MDV-Vorschlag ist. Dieser dient als Input für die Festlegung, deren Ergebnis, der \glqq MDV inklusive \acs{CoP}-Puffer \grqq , als Zwischenereignis dargestellt wird. Nachdem der AK Zertifizierung und die Projektleitersitzung informiert wurden, prüft ein Gateway die Bestätigung, andernfalls führt der Ablauf zur Festlegung zurück.
\newline
Dementsprechend wird in der dritten Phase der finale MDV nach demselben Muster unter Einbezug der Zertifizierungsergebnisse ermittelt und festgelegt, in den Fachgremien beschlossen und bei fehlendem Beschluss erneut festgelegt. Nach der Freigabe teilt ein UND-Gateway den Ablauf in die Eintragung im \acs{AZST} und die Übermittlung an die Abteilung \glqq MO/QZE\grqq , bevor der Prozess mit dem Endereignis endet.
\newline
Um den in Abschnitt \ref{sec:erstaufnahme} festgelegten Detaillierungsgrad einzuhalten, wurde jede Aktivität mit einer Tätigkeitsbeschreibung hinterlegt. Am 23. März 2026 wurde das Prozessmodell mit den direkten Prozessbeteiligten abgestimmt und fachlich bestätigt. Damit sind die in Kapitel \ref{chap:problemkontext} angenommenen Ablaufschritte sowie die Verantwortlichkeiten erstmals in einem Prozessmodell dargestellt. Dies bildet die Grundlage für die Verortung in Abschnitt \ref{sec:verortung}.

\section{Einordnung des MDV-Modells in die Prozessarchitektur}
\label{sec:verortung}

Zunächst wurde die Ebene des MDV-Prozesses in der Prozessarchitektur bestimmt und anschließend geprüft, ob er einer der Kategorien Management-, Kern- oder Unterstützungsprozesse zugeordnet werden kann. Die Prozessarchitektur von Mercedes-AMG umfasst vier Ebenen. Der MDV-Prozess liegt auf der vierten und damit untersten Ebene \citeintern{signavio_prozesslandkarte2026}.
\newline
Die Kategorien Management-, Kern- und Unterstützungsprozesse werden in der Literatur auf der obersten Ebene der Prozessarchitektur, der Prozesslandkarte, verwendet (vgl. Abschnitt~\ref{sec:Prozesslandkarte}). Da der MDV-Prozess auf der vierten Ebene liegt, wurde er keiner dieser Kategorien eigenständig zugeordnet. Seine Einordnung erfolgt stattdessen über die Hierarchie, d. h. über den übergeordneten \acs{WLTP}-Prozess, aus dem er als Teilprozess hervorgeht  \citeintern{wltp_eu7_workshop2026}, \citeintern{signavio_prozesslandkarte2026}. Dies ist mit der Literatur vereinbar, da die Bezeichnungen nicht einheitlich definiert sind und die Einordnung je nach Organisation abweichen kann \cite{ottenlips2025geschaftsprozessmanagement}.
\newline
Die Einordnung ist vorläufig, da der WLTP-Gesamtprozess derzeit an die Abgasnorm Euro 7 angepasst wird und die Veröffentlichung in Signavio erst nach der Gremienfreigabe vorgesehen ist \citeintern{wltp_eu7_workshop2026}. Dementsprechend ist die Einordnung bei Änderungen des Gesamtprozesses erneut zu prüfen. 




% Zunächst wurde der MDV-Prozess einer Kategorie zugeordnet, anschließend seine Ebene bestimmt und abschließend seine Beziehung zu angrenzenden Prozessen betrachtet. Der MDV-Prozess wurde den Kernprozessen zugeordnet, da er sich auf das Produkt bezieht und den deklarierten Wert liefert, der für die Typgenehmigung und die Kundeninformation benötigt wird (vgl. Abschnitt~\ref{sec:Prozesslandkarte}). Eine Zuordnung zu den Unterstützungsprozessen wurde verworfen, da der Prozess den übrigen Prozessen keine Ressourcen oder Grundlagen bereitstellt, sondern einen produktbezogenen Wert liefert. Eine Zuordnung zu den Managementprozessen wurde ebenfalls verworfen, da der Prozess keine Vorgaben für andere Prozesse festlegt, sondern selbst einen konkreten Wert ermittelt. Hinsichtlich der Ebene ist der MDV-Prozess ein Teilprozess des \acs{WLTP}-Prozesses \citeintern{wltp_eu7_workshop2026}. Da seine 13 Prozessschritte in BPMN 2.0 modelliert wurden, liegt er auf der untersten Ebene der Prozessarchitektur, auf der der Ablauf eines einzelnen Prozesses im Detail abgebildet wird (vgl. Abschnitt~\ref{sec:Prozesslandkarte}). Entsprechend den Beziehungsarten aus Abschnitt~\ref{sec:Prozesslandkarte} geht der MDV-Prozess durch Zerlegung als Teilprozess aus dem \acs{WLTP}-Prozess hervor. Zu angrenzenden Prozessen steht er in einer Sequenzbeziehung, da Prozesse, die die Fahrzeuge bereitstellen und Entwicklungs-, Simulations- und Vortestergebnisse liefern, vorgelagert sind und ein Prozess, der den ermittelten MDV weiterverwendet, nachgelagert ist \citeintern{wltp_eu7_workshop2026}. Daneben bestehen durch Spezialisierung marktspezifische Varianten des Prozesses, von denen in der vorliegenden Arbeit nur der Markt Westeuropa modelliert wurde. 